\section{Conclusion}
\label{sec:conclusion}

We propose an active GLA architecture for a two-dimensional airfoil with a hinged
trailing-edge flap: the aerodynamic loads are supplied by an LDNet, which acts as the internal
model of a preview MPC fed by fused lidar measurements. The LDNet is trained on full-order FSI
co-simulations instead of being derived from a simplified aerodynamic theory, so that it
reproduces the viscous, unsteady and nonlinear response to gust and flap while remaining cheap
enough to be evaluated at every step of the receding horizon.

The training set is a campaign of $146$ partitioned co-simulations spanning gust-only,
flap-only and combined excitations. Training minimises a rollout loss, in which the LDNet is
composed with the structural integrator and advanced autonomously over hundreds of steps, so
that the quantity minimised is the one the controller depends on at run time. The training
regime, rather than the capacity of the LDNet, is what decides whether it can be used in closed
loop: teacher forcing produces LDNets that are accurate in one-step replay but unstable and
flap-blind in autonomous rollout, whereas rollout training restores both the stability and the
sensitivity of the predicted loads to the flap, and a damped latent update encodes the stable
aeroelastic equilibrium that the undamped LDNet lacks.

The retained LDNet, with ten latent states, the full six-signal input and six hidden layers,
reconstructs the sectional loads on the held-out test set with a combined $\mathrm{NRMSE}$ of
$1.3\cdot10^{-3}$; queried away from the airfoil surface, the same reconstruction network
returns the distributed velocity field, a capability that the controller does not use and that
the meshless formulation provides at no additional training cost. In closed loop, the
architecture reduces the peak lift excursion by more than $81\%$ in the central portion of the
gradient range and by $39$--$95\%$ over the eighteen cases of the CS-25.341 discrete-gust
envelope, without actuator saturation and without violating the structural stability check. The
degradations at the two corners of the envelope are explained by physical limits, the actuator
rate at the sharp--severe corner and the pitch trade-off at the long-gust corner, rather than
by the LDNet.

The sensing requirements are mild. The fusion layer keeps the mean reduction within about one
percentage point of the noise-free value up to per-shot errors of $10\%\,W_0$, and within $1.5$
percentage points under a range-proportional lidar model whose raw error reaches $10$--$30\%$
of $W_0$. The systematic errors that do break the loop are asymmetric in a favourable
direction: the controller absorbs an under-read gust down to at least $-10\%\,W_0$ with no
loss and tolerates a preview time shift of $\pm10$~ms symmetrically with no flag, while an
over-read gust degrades the reduction gradually from $+2\%\,W_0$ and only destabilises the
loop at $+10\%\,W_0$.

We evaluate the closed loop on the reduced aeroelastic system driven by the true gust. We
also evaluate the same controller, still predicting with the LDNet surrogate, on the
full-order model in place of the reduced one. This full-order closed loop confirms the alleviation, with a peak-excursion
reduction of $80.7\%$ on Test~1 and $74.6\%$ on Test~2, in reasonable agreement with the
reduced-system values above.

Some aspects, where there is possibility for improvement of the proposed architecture, have not
been fully explored in this work. The aeroelastic system is a two-degree-of-freedom
airfoil with a single trailing-edge flap, and the transfer of the methodology to a flexible
three-dimensional wing involves structural modes and several control surfaces that the present
formulation does not address. The depth ladder ends at $L=24$, where a plain $\tanh$ stack
attenuates the gradient layer after layer; residual connections would extend it. The design of
the simulation campaign was not optimised either, and automatic procedures to select the
training trajectories would reduce the number of full-order simulations required.

Finally, prior knowledge about the aeroelastic system already enters the LDNet through its
formulation rather than through its training data: the damped latent update encodes the stable
equilibrium towards which the airfoil relaxes, and the tail-compressing output layer biases the
reconstruction towards the load extremes that size the structure. The same principle suggests a
second-order latent formulation: letting the dynamics network predict the second derivative of
the latent state rather than the first would give the latent dynamics the mass--spring--damper
structure of the structural equation of motion, which the present first-order recurrence has to
learn from data instead.